
\section*{8. Solvable or nilpotent group schemes}
\phantomsection\label{sec:sga3-VIB-8}%
\addcontentsline{toc}{section}{8. Solvable or nilpotent group schemes}

\medskip
\phantomsection\label{par:sga3-VIB-8.1}%
\noindent\textbf{8.1.---}
Let \(\mathcal C\) be a category equipped with a topology \(T\)
(cf.\ Expos\'e~IV, \S4). For every presheaf \(P\) on
\(\mathcal C\), write \(P^\flat\) for the associated sheaf.

Let \(G\) be a presheaf of groups on \(\mathcal C\), let \(A\)
and \(B\) be two group subpresheaves of \(G\), and let
\(\operatorname{Comm}(A,B)\) be the group presheaf of
commutators of \(A\) and \(B\) in \(G\). Thus, for every
\(S\in\operatorname{Ob}\mathcal C\),
\(\operatorname{Comm}(A,B)(S)\) is the subgroup of \(G(S)\)
generated by the commutators
\[
  aba^{-1}b^{-1},
  \qquad a\in A(S),\quad b\in B(S).
\]
We write
\phantomsection\label{form:sga3-VIB-8.1-CommT}%
\[
  \operatorname{Comm}_T(A,B)
    =\operatorname{Comm}(A,B)^\flat.
\]

In the proof of \hyperref[prop:sga3-VIB-8.2]{8.2}, we shall
need the following results.%
\footnote{\label{note:sga3-VIB-81}Editor's note: these results
were stated without proof in the original. We have highlighted
them as \hyperref[lem:sga3-VIB-8.1.1]{Lemmas 8.1.1} and \SGARefLink{lem:sga3-VIB-8.1.2}{8.1.2} and supplied the details of the
proofs.}

\medskip
\phantomsection\label{lem:sga3-VIB-8.1.1}%
\noindent\textbf{Lemma 8.1.1.---}
\emph{Let \(A\subset G\) be sheaves of groups, with \(A\)
invariant in \(G\).}
\begin{enumerate}
\item[\textup{\SGAInlineRefTarget{lem:sga3-VIB-8.1.1:item:i}(i)}]\phantomsection\label{concl:sga3-VIB-8.1.1-i}
  \emph{\(\operatorname{Comm}_T(G,A)\) is the smallest
  invariant group subsheaf \(C\) of \(G\) such that the sheaf
  \((A/C)^\flat\), associated with the quotient presheaf
  \(A/C\), is central in \((G/C)^\flat\).}

\item[\textup{(ii)}]\phantomsection\label{concl:sga3-VIB-8.1.1-ii}
  \emph{In particular, \(\operatorname{Comm}_T(G,G)\) is the
  smallest invariant group subsheaf \(C\) of \(G\) such that
  the quotient sheaf \((G/C)^\flat\) is commutative.}
\end{enumerate}

Clearly, \textup{\SGARefLink{concl:sga3-VIB-8.1.1-ii}{(ii)}} is the special case \(A=G\) of
\textup{\SGARefLink{concl:sga3-VIB-8.1.1-i}{(i)}}, so it is enough to prove \textup{\SGARefLink{concl:sga3-VIB-8.1.1-i}{(i)}}. Let \(C\)
be a group subsheaf of \(A\), invariant in \(G\), such that the
quotient sheaf \((A/C)^\flat\) is central in
\((G/C)^\flat\). By Lemma IV, \SGARefLink{sga3:iv:lemma:4:4:8:1}{4.4.8.1}, the presheaves \(A/C\)
and \(G/C\) are separated. Hence, by IV, \SGARefLink{sga3:iv:proposition:4:3:11}{4.3.11}, all the
morphisms in the following diagram are monomorphisms:
\begin{center}
  \phantomsection\label{diag:sga3-VIB-8.1.1-monomorphisms}%
  \begin{tikzcd}[row sep=large,column sep=large]
    A/C \arrow[r,hook] \arrow[d,hook] &
    G/C \arrow[d,hook] \\
    (A/C)^\flat \arrow[r,hook] &
    (G/C)^\flat
  \end{tikzcd}
\end{center}
Since \((A/C)^\flat\) is central in \((G/C)^\flat\), the
presheaf \(A/C\) is central in \(G/C\). It follows that
\[
  \operatorname{Comm}(G,A)\subset C,
\]
and hence, by IV, \SGARefLink{sga3:iv:corollary:4:3:12}{4.3.12}, \(C\) contains
\(\operatorname{Comm}_T(G,A)\).

Conversely, \(\operatorname{Comm}(G,A)\) is an invariant,
separated group subpresheaf of \(A\)
(cf.\ IV, \SGARefLink{sga3:iv:definition:4:3:1}{4.3.1}). Therefore, by IV,
\SGARefLink{sga3:iv:lemma:4:4:8:2}{4.4.8.2}\textup{\SGARefLink{lem:sga3-VIB-8.1.1:item:i}{(i)}} and IV, \SGARefLink{sga3:iv:proposition:4:3:11}{4.3.11},
\[
  C=\operatorname{Comm}_T(G,A)
\]
is an invariant group subsheaf of \(A\) containing
\(\operatorname{Comm}(G,A)\). Consequently, the presheaf
\(A/C\) is central in \(G/C\), and hence, by IV,
\SGARefLink{sga3:iv:lemma:4:4:8:2}{4.4.8.2}\textup{\SGARefLink{sga3:iv:lemma:4:4:8:2:item:ii:01}{(ii)}}, \((A/C)^\flat\) is central in
\((G/C)^\flat\). This proves \hyperref[lem:sga3-VIB-8.1.1]{Lemma 8.1.1}.
